% !TeX spellcheck = de_DE

Die Umsetzung des PoC orientiert sich an der in Kapitel \ref{ch:konzeption} vorgestellten Konzeption und besteht aus zwei zentralen Komponenten.
Die erste Komponente bekommt die berechneten Metriken des Textes von einem weiteren Service. 
Anhand von gesetzten Anforderungen an die Metriken, führt die Komponente die Bewertung der Qualitätsanforderungen durch.
Zur Berechnung der Qualitätsmetriken wird ein bestehender Metrics Service verwendet.
Dieser stellt verschiedene Kennzahlen zur Bewertung von LS+-Texten bereit und bildet die Grundlage der Analyse (Siehe Kapitel \ref{ch:metrikenanalyse}).

Die zweite Komponente steuert den eigentlichen Verbesserungsprozess, dabei wird die Bewertung der ersten Komponente abgerufen und je nach Ergebnis eine weitere Iteration initialisiert oder abgebrochen.
Wenn die Qualitätsanforderungen nicht erfüllt sind, wird ein Feedback erzeugt und für die nächste Textgenerierung verwendet.
Der resultierende Text wird anschließend erneut bewertet, wodurch ein geschlossener Feedbackkreislauf entsteht.